%% kapa: exam=1 round=25 year=2014 period=1 form=A subject=civil_law subject_ko=민법 seq=31 number=31
%% source: Q-Net 감정평가사 기출문제 (cst00309) artlSeq=5101207 file=extracted/1008/5101207/2014년 제25회 감정평가사 1차시험 기출문제/1교시(A형).hwp
%% conversion: kordoc 4.24.51 -> tools/gen1.py (auto); review=proofread-sonnet; flags=none
\begin{question}{31}
민법상 부동산의 공동소유에 관한 설명으로 옳지 않은 것은? (다툼이 있으면 판례에 의함)
\begin{choices}
  \item 공유자는 자기지분을 원칙적으로 자유롭게 처분할 수 있으나, 합유지분의 처분에는 제한이 있다.
  \item 합유물의 분할청구는 인정되지 않지만, 조합체의 해산 시에는 분할이 인정된다.
  \item 공유물의 관리에 관한 사항은 특약이 없는 한, 지분의 과반수로 결정한다.
  \item 비법인사단의 사원은 정관 기타 계약에 정함이 없는 한, 그 지위를 상실하면 총유물에 관한 권리도 소멸한다.
  \item 비법인사단인 교회의 재산에 대하여 각 교인은 그 지분권을 인정받을 수 있다.
\end{choices}
\end{question}
\answer{5}
